\documentclass[ ../main.tex]{subfiles}
\providecommand{\mainx}{..}
\begin{document}
\section{Data types with invariants}
\label{sec:invariants}

The type constructors of the previous sections (sum, product, exponential) produce \emph{free} types: every element of the carrier set is a valid value.
In practice, data types carry \emph{invariants}, predicates that constrain which values are admissible.
This section formalizes how approximation interacts with such constraints.

\subsection{Definition}

\begin{definition}[Data type with invariant]
\label{def:data_type_invariant}
A \emph{data type with invariant} is a pair $(\Set{T}, \Fun{P})$ where $\Set{T}$ is a type (a finite set) and $\Fun{P} \colon \Set{T} \to \Bool$ is a predicate.
The set of \emph{valid values} is
\begin{equation}
\Set{T}_{\Fun{P}} = \SetBuilder*{t \in \Set{T}}{\Fun{P}(t) = \True}\,,
\end{equation}
and the set of \emph{invalid values} is its complement $\Set{T} \setminus \Set{T}_{\Fun{P}}$.
\end{definition}

A data type with invariant is equivalently a \emph{subtype}: a subset of the carrier set selected by the predicate.
In relational terms, when $\Set{T}$ is a product type $\Set{A}_1 \times \cdots \times \Set{A}_k$, the invariant defines a \emph{relation} (a subset of the Cartesian product), and the valid values are the tuples in that relation.

\subsection{Invariant violation under approximation}

When a value $x^* \in \Set{T}_{\Fun{P}}$ is approximated, the approximate value $\tilde{x} \in \AT{\Set{T}}$ is drawn from the full type $\Set{T}$, not the constrained subtype $\Set{T}_{\Fun{P}}$.
The approximation mechanism has no knowledge of the invariant: it can produce any value in $\Set{T}$ according to the confusion matrix $M$.

\begin{definition}[Invariant violation probability]
\label{def:violation_prob}
Given a data type $(\Set{T}, \Fun{P})$ and an approximate value $\tilde{x} \in \AT{\Set{T}}$ with latent value $x^* \in \Set{T}_{\Fun{P}}$, the \emph{invariant violation probability} is
\begin{equation}
\label{eq:violation_prob}
\Prob{\Fun{P}(\tilde{x}) = \False \Given x^* = s} = \sum_{\substack{t \in \Set{T} \\ \Fun{P}(t) = \False}} M_{s, t}\,.
\end{equation}
\end{definition}

This is the total probability mass that the confusion matrix places on invalid values, given the latent value $s$.
It depends on both the error model (through $M$) and the structure of the invariant (through the set $\Set{T} \setminus \Set{T}_{\Fun{P}}$).

\begin{theorem}[Invariant violation bound]
\label{thm:violation_bound}
Let $(\Set{T}, \Fun{P})$ be a data type with invariant.
Let $\tilde{x} \in \AT{\Set{T}}$ be an approximate value with latent value $x^* \in \Set{T}_{\Fun{P}}$.
If the confusion matrix $M$ is a first-order (symmetric) channel with error rate $\fprate$, uniformly distributing error mass over all $\Card{\Set{T}} - 1$ incorrect values, then
\begin{equation}
\label{eq:violation_bound}
\Prob{\Fun{P}(\tilde{x}) = \False} = \fprate \cdot \frac{\Card{\Set{T} \setminus \Set{T}_{\Fun{P}}}}{\Card{\Set{T}} - 1}\,.
\end{equation}
\end{theorem}
\begin{proof}
Under a first-order symmetric channel, each incorrect value receives probability mass $\fprate / (\Card{\Set{T}} - 1)$.
Summing over the $\Card{\Set{T} \setminus \Set{T}_{\Fun{P}}}$ invalid values gives the result.
\end{proof}

\begin{corollary}
If the invariant excludes only a small fraction of $\Set{T}$, the violation probability is correspondingly small.
In the limiting case where $\Card{\Set{T}_{\Fun{P}}} = \Card{\Set{T}}$ (no constraint), the violation probability is zero.
At the opposite extreme, if $\Card{\Set{T}_{\Fun{P}}} = 1$ (a single valid value), then $\Prob{\Fun{P}(\tilde{x}) = \False} = \fprate$.
\end{corollary}

\subsection{Product types with invariants}

When the type is a product $\Set{A} \times \Set{B}$ with component-wise independent approximation, the invariant violation probability can be expressed in terms of the component error rates.

Let $\Fun{P} \colon \Set{A} \times \Set{B} \to \Bool$ be an invariant on the product type.
If the approximate value $(\tilde{a}, \tilde{b}) \in \AT{\Set{A}} \times \AT{\Set{B}}$ has independent component confusion matrices $M_{\Set{A}}$ and $M_{\Set{B}}$, then
\begin{equation}
\label{eq:product_violation}
\Prob{\Fun{P}(\tilde{a}, \tilde{b}) = \False \Given (a^*, b^*)} = \sum_{\substack{(a, b) \in \Set{A} \times \Set{B} \\ \Fun{P}(a, b) = \False}} (M_{\Set{A}})_{a^*, a} \cdot (M_{\Set{B}})_{b^*, b}\,.
\end{equation}

The factored form of the joint probabilities (from the Kronecker structure of product types) makes this sum tractable.
However, the invariant $\Fun{P}$ may couple the components, so the sum does not factor further in general.

\subsection{Worked example: approximate rational numbers}
\label{sec:approx_rational}

\begin{example}[Approximate rational number]
Consider the type $\mathrm{Rational} = (\IntSet, \IntSet)$ with invariant $\Fun{P}(n, d) = (d \neq 0)$.
In a finite implementation with $\IntSet = \{-K, \ldots, K\}$ for some bound $K$, the type has $\Card{\IntSet}^2 = (2K + 1)^2$ values, of which $(2K + 1)$ have $d = 0$ (one for each numerator).

Under independent, first-order symmetric approximation of each component with error rate $\fprate$, the probability that the denominator becomes exactly $0$ when the true denominator is $d^* \neq 0$ is
\begin{equation}
\Prob{\tilde{d} = 0 \Given d^* \neq 0} = \frac{\fprate}{2K}\,,
\end{equation}
since there are $2K$ incorrect denominator values and the error mass $\fprate$ is distributed uniformly among them.

For a 32-bit integer ($K = 2^{31} - 1$), this probability is approximately $\fprate / (2^{32} - 2) \approx \fprate \cdot 2.3 \times 10^{-10}$, which is negligible for any reasonable error rate.
The invariant is ``robust'' in the sense that the invalid subspace ($d = 0$) is a negligible fraction of the full type.
\end{example}

More generally, a data type invariant is robust under approximation when the fraction $\Card{\Set{T} \setminus \Set{T}_{\Fun{P}}} / \Card{\Set{T}}$ is small.
Invariants that exclude a large fraction of the carrier set are fragile: approximation is likely to produce an invalid value.

\subsection{Detecting and handling violations}

When an approximate value may violate an invariant, there are several strategies.

\paragraph{Runtime validation.}
If the invariant $\Fun{P}$ is computable, the approximate value can be checked at runtime.
This converts the data type to an optional type: $\AT{\Set{T}_{\Fun{P}}} \hookrightarrow \AT{\Set{T}} \to \mathrm{Maybe}(\Set{T}_{\Fun{P}})$, where the result is $\mathrm{Nothing}$ when the invariant is violated.
The probability of obtaining $\mathrm{Nothing}$ is exactly the invariant violation probability from \cref{def:violation_prob}.

\paragraph{Projection to valid subspace.}
In some cases, there is a natural projection $\pi \colon \Set{T} \to \Set{T}_{\Fun{P}}$ that maps invalid values to nearby valid ones (e.g., clamping a denominator of $0$ to $1$).
This introduces a deterministic correction that modifies the effective confusion matrix but eliminates invariant violations entirely.

\paragraph{Accepting the violation.}
When the violation probability is negligible (as in the rational number example), the violation can simply be accepted as part of the error model.
The data type's operations must then be robust to occasional invalid inputs, or the violation probability must be accounted for in downstream error analysis.

\end{document}
